\chapter{Guía del Director de Juego}

\section{Tono de la aventura: Acero, nostalgia y misterio}

El encanto de \emph{Entre barcos y acero} radica en el contraste entre la vida cotidiana de un grupo de chavales de clase trabajadora en 1992 y la intrusión de un horror cósmico y conspiranoico en su propio patio trasero. 

\begin{itemize}
  \item \lead{Nostalgia sensorial:} Recuerda a los jugadores los olores (carbón mojado, fritanga de taberna, chicle de fresa), los sonidos (frenos de bicicleta chirriando, estática de radio, trenes de cercanías) y las limitaciones de la época: no hay internet, no hay teléfonos móviles y la noche pertenece a quienes conocen los atajos.
  \item \lead{Los adultos no lo entenderán:} La policía municipal está desbordada por las huelgas de AHV; los padres trabajan turnos dobles o temen perder el empleo; y el Dr. Kessler tiene compradas a las autoridades locales. Los chavales están solos frente al misterio.
\end{itemize}

\section{El reloj del Eclipse (Gestión del tiempo)}

El clímax de la aventura está fijado en la noche del \lead{8 de diciembre de 1992}, momento en que el eclipse total de luna multiplicará la potencia del Éter. Utiliza un \lead{Reloj de Tensión de 6 segmentos}:

\begin{CajaRegla}{Reloj de Tensión hacia el Eclipse}
Cada vez que el grupo descanse para recuperar fuerzas, falle estrepitosamente una tirada de infiltración o Alex sufra una sobrecarga descontrolada, avanza 1 segmento en el reloj:
\begin{enumerate}
  \item \textbf{Segmento 1:} El cielo nocturno adquiere tonos verdosos; las brújulas empiezan a desviarse.
  \item \textbf{Segmento 2:} Las emisiones de radio y televisión local (ETB) sufren interferencias; Paco Lobatón emite un especial sobre La Arboleda.
  \item \textbf{Segmento 3:} Patrullas de seguridad privada de B\&W montan controles clandestinos en las salidas de Sestao hacia La Reineta.
  \item \textbf{Segmento 4:} El agua de la ría del Nervión hierve en pequeñas burbujas bioluminiscentes; los perros del pueblo aúllan sin descanso.
  \item \textbf{Segmento 5:} El Dr. Kessler moviliza a todos sus operativos para cercar el refugio de los chavales.
  \item \textbf{Segmento 6 (Medianoche del 8 de diciembre):} El eclipse comienza. La luna se torna de color rojo sangre; el Éter resuena en toda la comarca.
\end{enumerate}
\end{CajaRegla}

\section{Matriz de pistas clave}

\begin{enumerate}
  \item \lead{Pista 1 (El informe de Francisco):} Archivo de 1974 en la Sociedad de Ciencias que detalla fósiles de amonites con aleaciones metálicas no terrestres.
  \item \lead{Pista 2 (La libreta de B\&W):} Encontrada en el coche de un ingeniero; contiene cálculos orbitales que relacionan el eclipse lunar con el pozo de La Arboleda.
  \item \lead{Pista 3 (El colgante de Alex):} Un fragmento cristalizado de Éter que Alex lleva al cuello y que pulsa cuando se acerca a maquinaria de B\&W.
  \item \lead{Pista 4 (El plano de la nave fósil):} Escondido en la caja fuerte de Francisco; muestra que la galería minera es en realidad la cubierta superior de una colonia alienígena.
\end{enumerate}
